\markdownRendererDocumentBegin
\markdownRendererSectionBegin
\markdownRendererSectionBegin
\markdownRendererHeadingTwo{第1讲 运动的描述}\markdownRendererInterblockSeparator
{}\markdownRendererSectionBegin
\markdownRendererHeadingThree{机械运动}\markdownRendererInterblockSeparator
{}物体的空间位置随时间的变化\markdownRendererInterblockSeparator
{}
\markdownRendererSectionEnd \markdownRendererSectionBegin
\markdownRendererHeadingThree{质点}\markdownRendererInterblockSeparator
{}有质量的点用于替代物体\markdownRendererParagraphSeparator
{}物体的大小、形状和体积对所研究问题的影响可以忽略不计\markdownRendererInterblockSeparator
{}
\markdownRendererSectionEnd \markdownRendererSectionBegin
\markdownRendererHeadingThree{参考系}\markdownRendererInterblockSeparator
{}处理和使用参考系满足四个性质：\markdownRendererParagraphSeparator
{}标准性 作为参考系的物体被假定是不动的，并且其他的物体的运动与静止都需要相对于这个物体。\markdownRendererSoftLineBreak
{}任意性 参考系的选择是任意的。\markdownRendererSoftLineBreak
{}统一性 研究多个物体需要选同一个物体。\markdownRendererSoftLineBreak
{}差异性 同一个物体的运动，在不同参考系下一般是不同的。\markdownRendererInterblockSeparator
{}
\markdownRendererSectionEnd \markdownRendererSectionBegin
\markdownRendererHeadingThree{时间}\markdownRendererInterblockSeparator
{}时间间隔\markdownRendererParagraphSeparator
{}时刻\markdownRendererInterblockSeparator
{}
\markdownRendererSectionEnd \markdownRendererSectionBegin
\markdownRendererHeadingThree{位移和路程}\markdownRendererInterblockSeparator
{}路程 运动轨迹的长度\markdownRendererParagraphSeparator
{}位移 从初位置指向末位置的有向线段\markdownRendererParagraphSeparator
{}注意\markdownRendererSoftLineBreak
{}后续大家的学习中，位移都是相对于一个固定的初始位置的，末位置是任意时刻的位置，初位置就是运动开始计时的位置。\markdownRendererSoftLineBreak
{}路程是长度，是一个值，不是轨迹。\markdownRendererInterblockSeparator
{}
\markdownRendererSectionEnd 
\markdownRendererSectionEnd \markdownRendererSectionBegin
\markdownRendererHeadingTwo{第2讲 运动快慢的描述——速度}\markdownRendererInterblockSeparator
{}速度（瞬时速度）纯正的速度\markdownRendererParagraphSeparator
{}平均速度 位移除以时间\markdownRendererParagraphSeparator
{}瞬时速率 速度的大小\markdownRendererParagraphSeparator
{}平均速率 路程除以时间\markdownRendererParagraphSeparator
{}x-t图像\markdownRendererSoftLineBreak
{}直线斜率代表速度，曲线呢？切线斜率代表速度，斜率为负的意味着速度方向和拟定的正方向相反。\markdownRendererInterblockSeparator
{}\markdownRendererSectionBegin
\markdownRendererHeadingThree{第3讲 描述物体速度变化快慢的物理量——加速度}\markdownRendererInterblockSeparator
{}加速度（平均加速度）\markdownRendererParagraphSeparator
{}你知道吗，这个定义其实是平均加速度的定义式，我们的瞬时加速度定义式有两种，一个是和瞬时速度类似的，另一个是牛二的。\markdownRendererParagraphSeparator
{}加速度与速度变化率的方向相同，加速度和速度完全没关系。就好像有人问你在哪？你说我刚出门走了200米一样。\markdownRendererParagraphSeparator
{}a与v同向，加速运动\markdownRendererSoftLineBreak
{}a与v反向，减速运动到0，后反向加速运动。\markdownRendererSoftLineBreak
{}a越大，速度变化越快，a越小，速度变化越慢。\markdownRendererParagraphSeparator
{}匀变速直线运动为加速度恒定的直线运动，大小和方向都不变。\markdownRendererParagraphSeparator
{}匀变速直线运动v-t公式\markdownRendererSoftLineBreak
{}末速度等于初速度+at。\markdownRendererInterblockSeparator
{}
\markdownRendererSectionEnd \markdownRendererSectionBegin
\markdownRendererHeadingThree{第4讲 匀变速直线运动规律}\markdownRendererInterblockSeparator
{}v-t图像怎么用？\markdownRendererParagraphSeparator
{}v-t图像面积\markdownRendererParagraphSeparator
{}匀变速直线运动x-t公式\markdownRendererParagraphSeparator
{}匀变速直线运动x-v公式\markdownRendererInterblockSeparator
{}
\markdownRendererSectionEnd \markdownRendererSectionBegin
\markdownRendererHeadingThree{第5讲 自由落体运动}\markdownRendererInterblockSeparator
{}5个基本公式选择\markdownRendererParagraphSeparator
{}自由落体运动 为什么重力在两极最大，赤道最小呢？\markdownRendererParagraphSeparator
{}竖直上抛运动 和匀变速的关系，常用的公式。\markdownRendererInterblockSeparator
{}
\markdownRendererSectionEnd \markdownRendererSectionBegin
\markdownRendererHeadingThree{第6讲 运动学综合提升}\markdownRendererInterblockSeparator
{}公式汇总点评\markdownRendererParagraphSeparator
{}图像理解，x-t图像的理解，v-t图像的理解。\markdownRendererInterblockSeparator
{}
\markdownRendererSectionEnd \markdownRendererSectionBegin
\markdownRendererHeadingThree{第7讲 相互作用力}\markdownRendererInterblockSeparator
{}重力 mg，万有引力的分力，目前记住规律，有个印象\markdownRendererParagraphSeparator
{}弹力 有形变，想要回复形变的力叫弹力，弹力的方向比较难：\markdownRendererSoftLineBreak
{}绳子只有沿绳拉力，杆子可以像绳子一样提供拉力和推力或支持力，同时也可以提供转动的力矩，也就是初中的杠杆。\markdownRendererSoftLineBreak
{}弹簧可以沿绳方向可以提供拉力和推力或支持力。\markdownRendererSoftLineBreak
{}有接触面的可以提供垂直接触面的弹力，固定杆能够提供的弹力是完全未知方向的。\markdownRendererSoftLineBreak
{}举几个例子\markdownRendererInterblockSeparator
{}
\markdownRendererSectionEnd \markdownRendererSectionBegin
\markdownRendererHeadingThree{第8讲 摩擦力}\markdownRendererInterblockSeparator
{}静摩擦力 相互静止，必须要有弹力，最大静摩擦力和正压力成正比，静摩擦力本身和正压力无关。\markdownRendererParagraphSeparator
{}滑动摩擦力 相互运动时，必须要有弹力，和正压力成正比。\markdownRendererInterblockSeparator
{}
\markdownRendererSectionEnd \markdownRendererSectionBegin
\markdownRendererHeadingThree{第9讲 力的合成与分解}\markdownRendererInterblockSeparator
{}用一个力替代之前的力的作用效果，这个力就是其他力的合力。\markdownRendererSoftLineBreak
{}将一个力按照平行四边形或者三角形法则可以进行合成与分解。\markdownRendererSoftLineBreak
{}力的大小范围 差的绝对值 0 和 和\markdownRendererParagraphSeparator
{}正交分解\markdownRendererParagraphSeparator
{}水平竖直分解，沿斜面和垂直斜面分解，后面有加速度了就沿加速度和垂直加速度分解，等等，这是正交分解的一些妙用。\markdownRendererParagraphSeparator
{}三角函数常见的\markdownRendererParagraphSeparator
{}3 4 5\markdownRendererSoftLineBreak
{}5 12 13\markdownRendererInterblockSeparator
{}
\markdownRendererSectionEnd \markdownRendererSectionBegin
\markdownRendererHeadingThree{第10讲 共点力平衡}\markdownRendererInterblockSeparator
{}受力分析怎么做？\markdownRendererParagraphSeparator
{}重力先画，不一定要画在重心位置，或者是中心点，挑一个好看的地点更好，比如图中已经有很多直线了，选择那些直线的交点。这样可以让图像更清晰。\markdownRendererParagraphSeparator
{}弹力是一个难点，牢记我们之前的规律应用他们，\markdownRendererParagraphSeparator
{}最后再画摩擦力，因为有弹力才能有摩擦力。\markdownRendererParagraphSeparator
{}总结为1重2弹3摩擦。最后题目中还有别的外力再补充。\markdownRendererParagraphSeparator
{}二力平衡，一定是等大反向，作用于同一个物体。，方向是一条直线上。\markdownRendererParagraphSeparator
{}不共线的三力平衡一定是形成一个首尾相连的密闭三角形。\markdownRendererParagraphSeparator
{}多个力作用平衡时，任意一个力与其他力的合力是二力平衡。\markdownRendererParagraphSeparator
{}在上面的过程中，唯一需要注意的就是，多个力的合力不要和其他任何一个分力同时存在，一旦选择了合成，那么这些力就应该消失。\markdownRendererParagraphSeparator
{}正交分解处理平衡问题，只需要在两个垂直的方向上，合力为0就是平衡。一个方向平衡，那么就是那个方向上合力为0.\markdownRendererParagraphSeparator
{}总结一下：三个力及以上，我们很少使用正交分解方法，最多使用正弦定理。四个力则基本上都要用正交分解解决。\markdownRendererInterblockSeparator
{}
\markdownRendererSectionEnd \markdownRendererSectionBegin
\markdownRendererHeadingThree{第11讲 受力分析综合提升专题}\markdownRendererInterblockSeparator
{}弹力、摩擦力判断有无\markdownRendererParagraphSeparator
{}假设法，假设不存在或者假设存在，判断能够保持平衡，如果可以那么假设正确，如果无法，那么假设错误。同样的先假设弹力再假设摩擦力，因为先有弹力才有摩擦力。\markdownRendererParagraphSeparator
{}弹簧弹力的计算，和伸长量成正比，一定要记住，弹簧有一个位置是原长位置。这个位置相当重要，大家在后面选修和牛顿运动定律都会有弹簧，这个万恶的东西，比较难。\markdownRendererParagraphSeparator
{}摩擦力，里面的计算一定要注意，静摩擦力是大小未知的，和平衡状态有关，最大静摩擦力和正压力相关，一般而言题目会假定最大静摩擦力等于滑动摩擦力，但是不是一定的，一定一定要注意这个问题，但是有的时候题目确实说法不够严谨的，尤其是校级考试，那没办法，就认为这件事情是默认的。\markdownRendererInterblockSeparator
{}
\markdownRendererSectionEnd \markdownRendererSectionBegin
\markdownRendererHeadingThree{第12讲 牛顿运动定律}\markdownRendererInterblockSeparator
{}伽利略理想斜面实验\markdownRendererParagraphSeparator
{}三个实验推断过程\markdownRendererParagraphSeparator
{}牛顿第一运动定律（惯性定律）\markdownRendererSoftLineBreak
{}质量是惯性的量度，质量越大惯性就越大。\markdownRendererSoftLineBreak
{}惯性是物体改变运动状态难易程度的体现\markdownRendererSoftLineBreak
{}为什么我们之前质点，保留了质量呢？就是因为质量会影响惯性，而惯性会影响物体运动状态改变。所以运动学问题需要质量。\markdownRendererParagraphSeparator
{}牛顿第三运动定律（相互作用力）\markdownRendererParagraphSeparator
{}牛顿第二运动定律（非常非常重要的定律）\markdownRendererSoftLineBreak
{}成功将运动学和力学变成了一个统一的学科，力学，因为运动就是力的反应。力可以影响运动。\markdownRendererSoftLineBreak
{}F=ma里面的加速度是瞬时加速度，这就是我们以前提到过的，这才是加速度的定义式，能够决定任意时刻的加速度，并且只要物体的力改变了，物体的加速度就会改变。\markdownRendererSoftLineBreak
{}不同方向的力能得到对应方向的加速度。\markdownRendererParagraphSeparator
{}并且我们之前的受力平衡，其实你会发现牛二里面如果加速度是0，就是平衡问题。\markdownRendererParagraphSeparator
{}国际单位制，（量纲法）你知道吗？这可是一个高考考点，简单来说就是单位可以计算。\markdownRendererSoftLineBreak
{}做选择题至此大家已经收获了一个强大的方法，量纲法，选择题做之前先用量纲法试一试。\markdownRendererSoftLineBreak
{}这些单位里面记一下基本物理量就行，比如力学就是长度和时间与质量，电学需要电流就行，这个能够定义电场和磁场。热学需要物质的量和温度，压强用力学就行，光学里面加上一个发光强度。
\markdownRendererSectionEnd 
\markdownRendererSectionEnd 
\markdownRendererSectionEnd \markdownRendererDocumentEnd